\section{Proper learning without a dual-plane search}\label{sec:proper}
Every grid family contains the all-positive row: take intercept $b=2N^2$. Proper learning therefore does not require testing random lines when the total negative mass is small. The first substantial negative rectangle supplies the missing localization information.

\begin{samepage}
\begin{theorem}[Proper RECTIFY]\label{thm:proper}
Suppose $H$ contains the all-positive hypothesis and has no negative $2\times2$. Given a rectangle sampler of coverage $r$, it has a distribution-independent proper SQ learner of expected error at most $\epsilon$, using
\begin{equation}\label{eq:propercomplexity}
 m_{\rm proper}\le m_{\rm RECTIFY}(\epsilon/4)
 +\lceil\log_2|X|\rceil+\lceil\log_2K\rceil
 +\lceil8/\epsilon\rceil+4
\end{equation}
queries at the tolerance of RECTIFY with accuracy $\epsilon/4$. For the incidence flips,
\[
 m_{\rm proper}=O(\log M+\epsilon^{-1}+\log(1/\epsilon)),\qquad
 \tau=\Omega(\epsilon/\log(1/\epsilon)).
\]
The cap implementation is polynomial-time in the explicit table. Its finite-net deterministic variant is also proper and has the same asymptotic query bound, with constants depending on the template dimension.
\end{theorem}
\end{samepage}
\begin{center}\fbox{\begin{minipage}{.94\linewidth}
\textbf{Proper RECTIFY.}\par\smallskip
\begin{enumerate}
\item Run RECTIFY at accuracy $\epsilon/4$, stopping at the first negative peel. If none occurs before termination, return the all-positive class member.
\item In that negative rectangle, use its singleton point, or test its singleton row. Return a row whose loss answer is at most $\epsilon/2$. If rejected, binary-search the unique negative target point in the rectangle.
\item Query the negative tail away from that point. For a tail answer greater than $\epsilon/4$, binary-search the star using negative-union queries and return the target. Otherwise test up to $\lceil8/\epsilon\rceil+1$ star rows and return one with loss answer at most $\epsilon/2$.
\end{enumerate}
\end{minipage}}\end{center}
\begin{proof}
Set $\theta=\epsilon/4$ and run RECTIFY at accuracy $\theta$, stopping at its first proposed negative peel. Until then all assigned labels, and its fallback, are positive. If it stops normally without a negative peel, output the all-positive member of $H$. The error is at most $\theta/2$. Its budget-failure probability is at most $\theta/2$.

Suppose a negative rectangle $S\times T$ reaches a peeling answer. The acceptance calculation in Theorem~\ref{thm:main} and $u>\theta/8$ give
\[
 D(S)>5r\theta/128.
\]
For $P_\theta=1+\lceil4r^{-1}\ln(8/\theta)\rceil$ and $\tau=\theta/(24P_\theta)$,
\[
 \frac{D(S)}\tau>\frac{15}{16}rP_\theta
 >\frac{15}{4}\ln(8/\theta)>18,
\]
since $\theta<1/16$. A peeling answer implies $D(S\cap\{h^*=+1\})\le3\tau$. Hence the negative mass in $S$ exceeds $15\tau$.

If $S=\{p\}$, then $h^*(p)=-1$ and $p$ is known. Otherwise the negative rectangle has a single row $h_0$. Query its full loss, accepting it when the answer is at most $\epsilon/2$. An accepted candidate has true loss at most $\epsilon/2+\tau$. If it is rejected, $h_0\ne h^*$, since the target's answer would be at most $\tau$. The two rows share at most one negative point. Therefore the negative measure in $S$ is concentrated at a unique point $p$ and has mass greater than $15\tau$. Locate $p$ by binary search over an arbitrary ordering of $S$: for a half $B\subseteq S$, query $D(B\cap\{h^*=-1\})$. The answer is at most $\tau$ when $p\notin B$, and greater than $14\tau$ when $p\in B$. Threshold $4\tau$ safely determines the half. This takes at most $\lceil\log_2|X|\rceil$ queries.

Now let $V_p=\{h:h(p)=-1\}$. Its negative supports away from $p$ are pairwise disjoint; otherwise two rows share two negative points. Query the target's negative mass away from $p$,
\[
 \beta=D(\{x\ne p:h^*(x)=-1\}).
\]
If the answer exceeds $\epsilon/4$, then $\beta>\epsilon/4-\tau$. For any half $W\subseteq V_p$, define
\[
 Q_W(x,y)=\1_{\{y=-1\}}\1_{\{x\ne p\}}
 \1_{\{x\in\bigcup_{h\in W}\{h=-1\}\}}.
\]
Its expectation is exactly $\beta$ if $h^*\in W$ and exactly zero otherwise. Thus threshold $\epsilon/8$ distinguishes the two cases, since $\tau<\epsilon/16$. Binary search finds the target row in at most $\lceil\log_2K\rceil$ queries. If two indexed rows coincide and contain only $p$ as a negative point, this branch cannot occur; in the positive-tail branch the target row is unique.

If the answer for $\beta$ is at most $\epsilon/4$, then $\beta\le\epsilon/4+\tau$. Test up to
\[
 J=\lceil8/\epsilon\rceil+1
\]
distinct members of $V_p$, again accepting an estimated full loss at most $\epsilon/2$. If fewer than $J$ members exist, these tests include the target. Otherwise their negative masses outside $p$ sum to at most one, by disjointness. One tested row $h$ has such mass at most $1/J<\epsilon/8$. It agrees at $p$, and therefore
\[
 L_{D,h^*}(h)\le\beta+D(\{x\ne p:h(x)=-1\})
 <3\epsilon/8+\tau.
\]
Its answer is less than $\epsilon/2$, because $2\tau<\epsilon/8$. Some test consequently passes. Every accepted row has loss at most $\epsilon/2+\tau$.

All nonfailure outcomes are proper and have loss at most $\epsilon/2+\tau$. Any budget failure can output the all-positive row, with loss at most one. Expected loss is bounded by
\[
 \epsilon/2+\tau+\theta/2<\epsilon.
\]
The extra query counts give \eqref{eq:propercomplexity}. Using the deterministic net implementation removes the budget-failure event and leaves the same proper completion. Its computation consists of cap tests, row-label tests, and unions of explicitly listed row supports, all polynomial in the table size.
\end{proof}

The all-light case requires no aspect-ratio argument. At the first negative peel, either the singleton candidate is already good, or its intersection with the target contains one heavy negative atom. The other members of the star at that atom have disjoint negative supports. This is a proper-learning theorem, not an exponential proper--improper separation.
